\chapter{Torque and Rotational Dynamics}

\mathTopics{\pgRef{chap:vectors}}

\section{Rotational Kinematics}

Rotational motion is the motion of spinning objects.
Instead of measuring distance, we measure angle.

Similarly to the linear kinematic equations under constant linear
acceleration (\cpageref{eq:linear-kinematic}),
there are \textbf{rotational kinematic equations} under constant
\textbf{angular acceleration},
namely
\begin{subequations}
    \begin{gather}
        {
            \Delta{\theta} = \omega_0 t + \frac{1}{2} \alpha t^2
            \qquad \mathrm{or} \qquad
            \theta_\mathrm{f} = \theta_0 + \omega_0 t + \frac{1}{2} \alpha t^2
        } \\
        \omega_\mathrm{f} = \omega_0 + \alpha t \\
        \omega_\mathrm{f}^2 = \omega_0^2 + 2 \alpha \Delta{\theta} \\
        \Delta{\theta} = \frac{\omega_0 + \omega_\mathrm{f}}{2} t
    \end{gather}
\end{subequations}
where
\begin{itemize}
    \item \(\theta_0\) is the initial angle,
        \(\theta_\mathrm{f}\) is the final angle,
        and \(\Delta{\theta}\) is the change in angle;
    \item \(\omega_0\) is the initial \textbf{angular velocity},
        and \(\omega_\mathrm{f}\) is the final angular velocity;
    \item \(\alpha\) is the angular acceleration;
    \item \(t\) is the time elapsed.
\end{itemize}

The \textbf{arc length} \(s\) is given by
\begin{equation}
    s = r \theta,
\end{equation}
where \(\theta\) is the angle in radians.
It does not work in degrees.

The \textbf{tangential velocity} \(v_\mathrm{t}\) is given by
\begin{equation}
    v_\mathrm{t} = r \omega.
\end{equation}

The \textbf{tangential acceleration} \(a_\mathrm{t}\) is given by
\begin{equation}
    a_\mathrm{t} = r \alpha.
\end{equation}

\subsection{Moment of Inertia}

\textbf{Moment of inertia} or \textbf{rotational inertia} (\(I\)) is
the angular analogue of mass,
describing how much a body resists rotational acceleration.

For discrete masses, the rotational inertia \(I\) is
\begin{equation}
    I = \sum_i m_i r_i^2,
\end{equation}
where \(m_i\) is the mass of each particle,
and \(r_i\) is the distance of each particle from the \textbf{axis of rotation}.

For continuous masses, the rotational inertia \(I\) is
\begin{gather}
    I = \int r^2 \odif{m} \\
    \begin{aligned}
        \lambda(r) &= \odv{m}{r} \\
        \odif{m} &= \lambda(r) \odif{r}
    \end{aligned} \\
    \therefore I = \int r^2 \lambda(r) \odif{r}.
\end{gather}

Imagine a cylindrical rod of uniform density \(\lambda\), length
\(\ell\), and mass \(M\)
rotating about an axis perpendicular to its length and passing
through the center.

\begin{figure}[ht]
    \centering
    \begin{tikzpicture}
    \draw (-6, -0.3) rectangle (6, 0.3);
    \draw[{|[scale=2]}-{|[scale=2]}] (-6, -0.667) -- (6, -0.667)
    node[midway, below] {\(\ell\)};
    \draw[dotted, line width = 0.333mm, line cap=round] (0, -0.5) -- (0, 1);
    \draw[->, thick] (0, 0) arc [
        start angle=-90,
        end angle=240,
        x radius=6mm,
        y radius=3mm
    ];
\end{tikzpicture}

\end{figure}

The moment of inertia can be derived as
\begin{equation} \label{eq:uniform-rod-rotational-inertia}
    \begin{aligned}
        I &= 2 \int_{0}^{\frac{\ell}{2}} r^2 \odif{m} \\
        &= 2 \int_{0}^{\frac{\ell}{2}} r^2 \lambda \odif{r} \\
        &= 2 \lambda \int_{0}^{\frac{\ell}{2}} r^2 \odif{r} \\
        &= 2 \lambda \left[\frac{1}{3} r^3 \right]^\frac{\ell}{2}_0 \\
        &= \frac{2 \lambda}{3} \frac{\ell^3}{8} \\
        &= \lambda \frac{\ell^3}{12} \\
        &= \frac{1}{12} M \ell^2
    \end{aligned}.
\end{equation}

Note that the \(r\) in these formulas means that objects have lower
moment of inertia
when the axis of rotation is closer to their center of mass
than when it is far from the center of mass,
e.g., it is easier to spin a stick at the middle than at either end.
A more interesting example is that a full can rolls down a slope
faster than an empty can does.
The full can's mass is more concentrated in the center,
reducing its moment of inertia and thus giving it greater acceleration.

If we already know the rotational inertia of an object
when rotating about an axis passing through its center of mass,
the \textbf{parallel axis theorem} allows us to determine the rotational inertia
when it is rotating around a new axis parallel to the original axis.
It states
\begin{equation}
    I = I_\mathrm{c.m.} + m d^2,
\end{equation}
where
\begin{itemize}
    \item \(I\) is the new moment of inertia.
    \item \(I_\mathrm{c.m.}\) is the moment of inertia
        around an axis passing through the center of mass.
    \item \(m\) is the mass of the object.
    \item \(d\) is the distance between the two parallel axes.
\end{itemize}

Using the uniform rod from \pgRef{eq:uniform-rod-rotational-inertia}
as an example again,
if we rotated the rod around its end instead of its center, its
rotational inertia would be
\begin{equation}
    \begin{aligned}
        I &= \frac{1}{12} M \ell^2 + M \ab(\frac{\ell}{2})^2 \\
        &= \frac{1}{3} M \ell^2.
    \end{aligned}
\end{equation}

\begin{figure}[ht]
    \centering
    \begin{tikzpicture}
    \draw (-6, -0.3) rectangle (6, 0.3);
    \draw[{|[scale=2]}-{|[scale=2]}] (-6, -0.667) -- (6, -0.667) node[midway, below] {\(\ell\)};
    \node [circle, fill=black, text=black, minimum size=0.1mm] (cm) {};
    \node [right=0.5mm of cm] {c.m.};
    \draw[dotted, line width = 0.333mm, line cap=round] (0, -0.5) -- (0, 1);
    \draw[dotted, line width = 0.333mm, line cap=round] (6, -0.5) -- (6, 1);
    \draw[->, thick] (6, 0) arc [start angle=-90, end angle=225, radius=0.4cm];
\end{tikzpicture}
\end{figure}

\section{Torque}

\textbf{Torque} (\(\vecb{\tau}\)) is the rotational analogue to linear force
and is given by the equation
\begin{equation}
    \begin{aligned}
        \vecb{\tau} &= \vecb{r} \times \vecb{F} \\
        \tau &= r F \sin\theta,
    \end{aligned}
\end{equation}
where
\begin{itemize}
    \item \(\vecb{\tau}\) is the torque vector and \(\tau\) is its magnitude.
    \item \(\vecb{r}\) is the position vector from the point of rotation
        to the point where the force is applied
        and \(r\) is its magnitude.
    \item \(\vecb{F}\) is the force applied and \(F\) is its magnitude.
    \item \(\theta\) is the angle between the force and the position vector.
\end{itemize}

The relationship between the direction of the torque
and the direction of the applied force is given by the \textbf{right-hand rule}.
See \pgRef{chap:vectors} for details.

\textbf{Newton's second law for rotational movement} states that
\begin{equation} \label{eq:newtons-2nd-law-rotational}
    \begin{aligned}
        \vecb{\tau} &= I \vecb{\alpha} \\
        \tau &= I \alpha.
    \end{aligned}
\end{equation}

\subsection{Atwood Machines}

An \textbf{Atwood machine} is a pair of masses connected by a rope
and hung over a pulley.

Let
\begin{itemize}
    \item \(m_1\) and \(m_2\) be
        the masses on the left and right respectively.
    \item \(M\) be the mass of the pulley.
    \item \(F_1\) be the total force exerted on \(m_1\),
        and similarly for \(F_2\).
    \item \(F_\mathrm{T1}\) be the force of tension exerted on \(m_1\)
        and similarly for \(F_\mathrm{T2}\).
    \item \(R\) be the radius of the pulley.
    \item \(\alpha\) be the angular acceleration of the pulley.
    \item \(a\) be the linear acceleration that both masses experience.
\end{itemize}

From \pgRef{eq:newtons-2nd-law-rotational},
\begin{equation*}
    \begin{aligned}
        \tau &= I \alpha \\
        &= I \frac{a}{R} \\
    \end{aligned}
\end{equation*}
\begin{equation}
    \begin{aligned}
        F_1 &= F_\mathrm{T1} - m_1 g \\
        F_2 &= F_\mathrm{T2} - m_2 g \\
        F_\mathrm{T1} - m_1 g &= m_1 a \\
        F_\mathrm{T1} &= m_1 a + m_1 g \\
        F_\mathrm{T2} - m_2 g &= m_2 a \\
        F_\mathrm{T2} &= m_2 a + m_2 g \\
    \end{aligned}
\end{equation}
\begin{equation}
    \begin{aligned}
        R F_\mathrm{T2} - R F_\mathrm{T1} &= \frac{1}{2} M R a \\
        F_\mathrm{T2} - F_\mathrm{T1} &= \frac{1}{2} M a \\
        m_2 a + m_2 g - \ab(m_1 a + m_1 g) &= \frac{1}{2} M a \\
        m_2 g - m_1 g &= \frac{1}{2} M a - m_2 a + m_1 a \\
        \ab(m_2 - m_1) g &= \ab(\frac{1}{2} M - m_2 + m_1) a \\
    \end{aligned}
\end{equation}
\begin{equation}
    \begin{aligned}
        a &= \frac{\ab(m_2 - m_1) g}{\frac{1}{2} M - m_2 + m_1} \\
        &= \frac{2 \ab(m_2 - m_1) g}{M - 2 \ab(m_2 + m_1)}
    \end{aligned}
\end{equation}

\section{Rolling}

Objects need friction to roll.

If there is no friction, the object simply slides rather than rolling,
which does not caused energy loss.
The tangential velocity is 0,
but the velocity of the center of mass is non-zero.

If the static friction is insufficient, the object \textbf{rolls with slipping},
which causes energy loss.
The tangential velocity is less than the velocity of the center of mass.

If the static friction reaches a critical point, it \textbf{rolls
without slipping},
which does not cause energy loss.
The tangential velocity and the velocity of the center of mass are equal, i.e.,
\begin{equation}
    a_\mathrm{t} = a_\mathrm{c.m.}
\end{equation}

Imagine an object rolling down a ramp without slipping.
Let
\begin{itemize}
    \item \(r\) be the radius of an object rolling down a ramp.
    \item \(\alpha\) be the angular acceleration of the rolling object.
    \item \(a\) be the linear acceleration.
    \item \(m\) be the mass of the rolling object.
    \item \(\theta\) be the angle of the ramp from the horizontal.
    \item \(\mu_\mathrm{s}\) be the static friction of the ramp.
    \item \(F_\mathrm{N}\) be the normal force experienced by the object.
    \item \(F_\mathrm{f}\) be the frictional force experienced by the object.
\end{itemize}

Deriving \(\mu_\mathrm{s}\):
\begin{equation}
    \begin{aligned}
        F_\mathrm{N} &= m g \cos\theta \\
        \therefore F_\mathrm{f} &= \mu_\mathrm{s} m g \cos\theta
    \end{aligned}
\end{equation}
\begin{equation} \label{eq:rolling-static-friction}
    \begin{aligned}
        m g \sin\theta - F_\mathrm{f} &= m a \\
        m g \sin\theta - \mu_\mathrm{s} m g \cos\theta &= m a \\
        \mu_\mathrm{s} &= - \frac{a}{g \cos\theta} + \tan\theta \\
    \end{aligned}
\end{equation}
\begin{equation} \label{eq:rolling-angular-accleration}
    \begin{aligned}
        \tau &= I \alpha \\
        r F_\mathrm{f} &= I \alpha \\
        \alpha &= \frac{r F_\mathrm{f}}{I} \\
        &= \frac{r \mu_\mathrm{s} m g \cos\theta}{I} \\
    \end{aligned}
\end{equation}
Plug in \pgRef{eq:rolling-angular-accleration}:
\begin{equation}
    \begin{aligned} \label{eq:rolling-accleration}
        a &= r \alpha \\
        &= \frac{r^2 \mu_\mathrm{s} m g \cos\theta}{I} \\
    \end{aligned}
\end{equation}
Plug \cref{eq:rolling-accleration} into \pgRef{eq:rolling-static-friction}.
\begin{equation}
    \begin{aligned}
        \mu_\mathrm{s} &=
        - \frac{
            \frac{r^2 \mu_\mathrm{s} m g \cos\theta}{I}
        }
        {g \cos\theta}
        + \tan\theta \\
        &= - \frac{\mu_\mathrm{s} m r^2}{I} + \tan\theta \\
        \mu_\mathrm{s} + \frac{\mu_\mathrm{s} m r^2}{I} &= \tan\theta \\
        \mu_\mathrm{s} \ab(1 + \frac{m r^2}{I}) &= \tan\theta \\
        \mu_\mathrm{s} &= \frac{\tan\theta}{1 + \frac{m r^2}{I}} \\
        &= \frac{I \tan\theta}{I + m r^2}
    \end{aligned}
\end{equation}
